\VerdierSemanticAnchor{front-title}
\VerdierPhysicalPageAnchor{0002}
\begin{titlepage}
  \centering
  {\Huge\bfseries 239\hfill ASTÉRISQUE\par}
  \vspace{1.5em}
  {\LARGE\bfseries 1996\par}
  \vfill
  {\Huge DERIVED CATEGORIES\par}
  \vspace{0.4em}
  {\Huge OF ABELIAN CATEGORIES\par}
  \vspace{3em}
  {\Large\bfseries Jean-Louis VERDIER\par}
  \vfill
  {\Large\bfseries SOCIÉTÉ MATHÉMATIQUE DE FRANCE\par}
  \smallskip
  {\small\bfseries Published with the support of the CENTRE NATIONAL DE LA RECHERCHE SCIENTIFIQUE\par}
\end{titlepage}

\VerdierSemanticAnchor{front-ams-classification}
\VerdierPhysicalPageAnchor{0003}
\chapter*{A.M.S. Subjects classification}
\addcontentsline{toc}{chapter}{A.M.S. Subjects classification}

14F20, 18E30, 18E35, 18G05, 18G10, 18G15, 18G20,18G35,
18G40, 55U15, 55U25, 55U30.

\clearpage
\VerdierSemanticAnchor{front-portrait}
\VerdierPhysicalPageAnchor{0004}
\thispagestyle{empty}
\vspace*{0.62\textheight}
\begin{center}
Jean-Louis VERDIER, 1977
\end{center}

\clearpage
\VerdierPhysicalPageAnchor{0005}
\thispagestyle{empty}
\null
\clearpage

\VerdierSemanticAnchor{front-preface}
\VerdierPhysicalPageAnchor{0006}
\chapter*{Preface}
\addcontentsline{toc}{chapter}{Preface}

Jean-Louis Verdier, who died tragically on 25 August 1989, defended his thesis
for the \emph{Doctorat d'État}, entitled \emph{Derived Categories of Abelian
Categories}, on 14 June 1967 in Paris. The present text is that of the
typescript of this thesis, in its final version. The editors confined
themselves to correcting minor errors. Only the part entitled
\emph{Introduction} had been deposited with the faculty. The body of the text
was unfinished at the time of the defense, and remained so.

The genesis of this work is fairly well known (\emph{cf.} [I]). We know that
Grothendieck conceived the idea of derived categories at the beginning of the
1960s in order to provide the framework of homological algebra required for
the vast generalization of the duality theorems that he had envisaged in the
wake of his address to the 1958 International Congress in Edinburgh. He had
proposed to Verdier, as a thesis subject, the construction of the formalism he
had envisaged. In 1963, Verdier brought out a fascicle of results summarizing
the essential content of the theory, \emph{Catégories dérivées, quelques
résultats (État 0)}, mimeographed by the IHÉS (reproduced in [SGA 4 1/2]). If
the notion of a derived category is due to Grothendieck, it was Verdier who had
earlier introduced that of a triangulated category. The axiomatic framework
presented there, inspired by stable homotopy theory (\emph{cf.} [P]), was to
prove surprisingly fruitful.

It was only several years later that Verdier undertook a comprehensive account
of his work. Why did he not finish it? Many hypotheses have been put forward.
I would incline toward the following one. The part that is missing is the
theory of derived functors. When Verdier reached that point in his writing,
derived categories had already been used extensively in algebraic geometry.
Various variants and generalizations were developing: the notion of pointwise
derivability à la Deligne, nonadditive derived functors (homotopical algebra à
la Quillen), filtered derived functors. From then on, the framework set by the
results fascicle proved inadequate. Finding the right point of view would have
required new foundational work which Verdier, interested at that time in
other problems, had, it seems to me, neither the leisure nor doubtless the
desire to carry out. As time passed, articles on derived categories appeared
and their use spread, the formalism invading other fields of mathematics
(such as linear partial differential equations). Completing his work then
became, in Verdier's eyes, an increasingly unnecessary task \ldots{}

\VerdierPhysicalPageAnchor{0007}
Incomplete though it is, Verdier's text nevertheless constitutes a valuable,
indeed irreplaceable, reference (for example, the original and elegant
presentation of spectral sequences given in it is, to my knowledge, found
nowhere else). One can only rejoice that it has finally seen the light of day.

\begin{flushright}
Paris, 27 April 1996\\
Luc Illusie
\end{flushright}

\section*{References}

\begin{description}
\item[{[I]}] L. Illusie, \emph{Catégories dérivées et dualité, travaux de
J.-L. Verdier}, L'Enseignement Mathématique, 36 (1990), 369-391. (Talk
delivered on 19 October 1989, during the ceremony in tribute to Jean-Louis
Verdier organized by Université Paris VII.)

\item[{[P]}] D. Puppe, \emph{On the formal structure of stable homotopy
theory}, Colloquium on Algebraic Topology, Aarhus Universitet (1962), 65-71.

\item[{[SGA 4 1/2]}] ``Cohomologie étale'', Séminaire de Géométrie Algébrique
du Bois-Marie, SGA 4 1/2, by P. Deligne, Lecture Notes in Mathematics 569,
Springer-Verlag (1977), 262-311.
\end{description}

